sections/03_framework.tex
\[ %%%%%%%%%%%%%%%%%%%%%%%%%%%%%%%%%%%%%%%%%%%%%%%%%%%%%%%%%%%%%%%%%%%%%%%%%%%%%%%% \section{PHYSGEN Framework} \label{sec:framework} %%%%%%%%%%%%%%%%%%%%%%%%%%%%%%%%%%%%%%%%%%%%%%%%%%%%%%%%%%%%%%%%%%%%%%%%%%%%%%%% The PHYSGEN framework introduces a physics-guided graph neural architecture for stable long-horizon prediction of nonlinear dynamical systems. The central idea behind PHYSGEN is to represent physical processes as evolving interaction graphs and incorporate physical knowledge directly into the computational mechanisms responsible for state representation, information propagation, temporal evolution, and stability regulation. Unlike conventional data-driven graph neural networks, where physical behavior is learned implicitly from simulation data, PHYSGEN treats physical principles as active components of the learning architecture. The framework combines four fundamental computational elements: \begin{enumerate} \item dynamic graph representation of physical systems; \item physics-guided message passing between interacting components; \item latent physical state evolution through differentiable dynamics; \item adaptive stability control for long-horizon prediction. \end{enumerate} The overall computational process can be expressed as a sequence of physics-guided transformations: \begin{equation} G_t \rightarrow H_t \rightarrow H_{t+1} \rightarrow G_{t+1}, \end{equation} where $G_t$ represents the physical system graph at time $t$, and $H_t$ denotes the latent physical state representation learned by the model. The proposed framework assumes that many complex physical systems can be represented as dynamic graphs: \begin{equation} G_t=(V_t,E_t,X_t), \end{equation} where: \begin{itemize} \item $V_t$ represents the set of physical entities or computational elements; \item $E_t$ represents interaction relationships between entities; \item $X_t$ contains time-dependent physical variables associated with graph nodes. \end{itemize} Depending on the application domain, graph nodes may correspond to finite-volume cells, mesh vertices, particles, molecular components, biological entities, or infrastructure elements. Graph edges encode local physical interactions, information exchange pathways, or coupling relationships between system components. The temporal evolution of the physical system is modeled as a graph dynamical process: \begin{equation} G_{t+1} = \Phi_{\theta}(G_t,\mathcal{P}), \end{equation} where $\Phi_{\theta}$ represents the learned evolution operator and $\mathcal{P}$ denotes embedded physical knowledge, including conservation principles, governing equations, constitutive relations, and domain-specific constraints. A key principle of PHYSGEN is that physical knowledge is integrated throughout the computational pipeline rather than applied exclusively during optimization. Therefore, the learned evolution operator is decomposed into several interacting components: \begin{equation} \Phi_{\theta} = \Phi_{stab} \circ \Phi_{phys} \circ \Phi_{msg} \circ \Phi_{enc}, \end{equation} where: \begin{itemize} \item $\Phi_{enc}$ is the physical state encoder; \item $\Phi_{msg}$ is the physics-guided message passing operator; \item $\Phi_{phys}$ represents physics-aware latent state evolution; \item $\Phi_{stab}$ performs adaptive stability regulation. \end{itemize} This decomposition reflects the main architectural hypothesis of PHYSGEN: \begin{quote} Reliable long-horizon prediction of physical systems requires physical guidance to influence not only the learning objective but also the internal mechanisms responsible for information transformation and temporal evolution. \end{quote} The proposed framework therefore moves beyond conventional surrogate modeling, where neural networks approximate numerical solvers independently of physical structure. Instead, PHYSGEN establishes a hybrid computational paradigm in which learning and physical reasoning are continuously coupled during prediction. Figure~\ref{fig:phygen_architecture} illustrates the overall architecture of PHYSGEN. The system receives an initial physical state, constructs a dynamic interaction graph, performs physics-guided message passing, evolves latent physical states, and generates future system configurations through stability-controlled recursive prediction. The framework is designed to be domain-independent and can be applied to a broad range of dynamical systems, including fluid flows, thermal processes, reaction--diffusion systems, mechanical structures, climate models, and digital twin environments. \] 
